% unit U6
\section{Scopes and Levels: The Compiler's Reading of \texorpdfstring{$\mathsf{any}$}{any}}\label{sec:levels}

% unit macros
\providecommand{\elej}[4]{\tgt{#1 \vdash_{\!e} #2 : #3 \le #4}}
\providecommand{\lamcc}[6]{\lambda^{#1}_{#2}(#3 : #4).\,#5\;\ann{#6}}
\providecommand{\frletex}[6]{\klet \langle #1, #2\rangle.\,#3\;\ann{#4;\,#5;\,#6}}
\providecommand{\appE}[2]{#1 \mathbin{\triangleleft} #2}
\providecommand{\ent}[1]{\mathit{ent}_{#1}}
\ifcsname remark\endcsname\else
  \theoremstyle{acmdefinition}\newtheorem{remark}[theorem]{Remark}
\fi

\mathtgt

\treeK reads $\cany$ by its position, the DOT way (Section~\ref{sec:source-any}). The Scala~3 compiler reads it by scope (Section~\ref{sec:cc-later}). A capability may flow into the local $\cany$ of its own scope or of a nested one, never into that of an enclosing one~\cite{scala3cc}. \treeCC models this reading. Every theorem here is in \treeCC, and its head says whether \treeCL has it too. The proof of Sections~\ref{sec:mini} and~\ref{sec:target} stays. What changes is what the roots of a capture set are: they now open each scope capability into the capabilities it stands for. Levels act while a program is typed. At run time every capability is at the outermost level, and a call turns the local $\cany$ of the body it enters into the universal capability.

\subsection{The escape}\label{sec:levels-program}

Section~\ref{sec:cc-later} quoted the escape of \lstinline{withFile}. The callback \lstinline|f => () => f| returns a closure of type \lstinline|() ->{f} File^{f}|, and the program asks for it at \lstinline|() => File^|, whose $\cany$ is bound outside \lstinline{withFile}. The reference refuses this by its rule on $\cfresh$ results, which Section~\ref{sec:levels-fresh} reads in the calculus. \treeK cannot pose the escape (Section~\ref{sec:source-any}). \treeCC can, since an arrow binds a capture variable and a parameter written $\cany$ reads as it. Its levels are meant to refuse it, since \lstinline{f} belongs to the callback's scope. The promise is that the file does not outlive the callback. What is proved is narrower. In a hand-made context for the callback's body, no evidence puts $\cset{f}$ below $\cset{\topc}$ (Section~\ref{sec:levels-inversion}), and no theorem rejects the escape program (Section~\ref{sec:levels-open}).

\subsection{Scopes, levels and the roots of a set}\label{sec:levels-scopes}

\begin{figure}[t]
\footnotesize
\color{tgtcol}
\[
\begin{array}{@{}l@{}}
\begin{array}{@{}l@{\quad}rcl@{}}
\text{Capture atoms} & c, r & ::= & x \mid \kappa \mid \sel{x}{\ell} \mid \topc \\
\text{Capture bounds} & B & ::= & \cstar \mid \scopeb \mid \instb{C} \qquad \text{in binders } \Gamma, \kappa : B \text{ and } \sigma, \kappa : B \\
\text{Shapes} & S & ::= & \bot \mid \sel{x}{\ell} \mid \pit[\kappa, x]{T}{E} \mid \obj[x]{\tel} \mid \boxt{T} \\
\text{Answers} & E & ::= & T \mid \exc{\kappa}{C}{T} \\
\text{Values} & v & ::= & \lamcc{A}{\varrho}{\kappa, x}{T}{t}{g} \mid \litc{A}{W}{W^{c}}{F} \mid \ev{box}\,a \mid \ev{cast}\,v\,d \\
\text{Terms} & t, u & ::= & \cdots \mid \ev{castE}\,t\,\eta \\
\text{Answer evidence} & \eta & ::= & \ev{plain}\,d \mid \ev{trans}\,\eta\,\eta \mid \cdots \\
\text{Capture evidence} & f, g & ::= & \cdots \mid \ev{level}\,c\,r \\
\text{Capture equality} & \psi & ::= & \cdots \mid \ev{instC}\,\kappa\,C \\
\text{Continuations} & K & ::= & \cdots \mid K \triangleright \ev{castE}\,\eta
\end{array} \\[6pt]
\begin{array}{@{}lcl@{}}
\sext{\sigma}{\sigma'} & \text{iff} & \sigma' = \sigma, \sigma'' \text{ where } \sigma'' \text{ binds no scope binder and no star binder} \\
\uses{\ev{castE}\,t\,\eta} & = & \uses{t} \\
\uses{K \triangleright \ev{castE}\,\eta} & = & \uses{K}
\end{array}
\end{array}
\]
\begin{mathpar}
\inferrule[C-Level]{r = \topc \text{ or } \Gamma \text{ binds } r : \scopeb \\ \lvlle{\Gamma}{c}{r}}{\clej{\Gamma}{\ev{level}\,c\,r}{\cset{c}}{\cset{r}}}
\and
\inferrule[CE-Inst]{\Gamma \text{ binds } \kappa \instb{C}}{\ceqj{\Gamma}{\ev{instC}\,\kappa\,C}{\cset{\kappa}}{C}}
\and
\inferrule[Sh-Pi]{\lej{\mathit{scope}(\Gamma)}{d}{T_2}{T_1} \\ \elej{\mathit{body}(\Gamma, T_2)}{\eta}{E_1}{E_2}}{\shpj{\Gamma}{\ev{pi}\,d\,\eta}{\pit[\kappa, x]{T_1}{E_1}}{\pit[\kappa, x]{T_2}{E_2}}}
\and
\inferrule[Ans-Plain]{\lej{\Gamma}{d}{T}{T'}}{\elej{\Gamma}{\ev{plain}\,d}{T}{T'}}
\and
\inferrule[Ans-Trans]{\elej{\Gamma}{\eta}{E_1}{E_2} \\ \elej{\Gamma}{\eta'}{E_2}{E_3}}{\elej{\Gamma}{\ev{trans}\,\eta\,\eta'}{E_1}{E_3}}
\and
\inferrule[T-App]{\atomj{\Gamma}{a}{\capt{(\pit[\kappa, x]{T}{E})}{C}} \\ \atomj{\Gamma}{b}{T[\kappa \mapsto \rt{b}]}}{\tmj{\Gamma}{a\,b}{E[\kappa \mapsto \rt{b}, x \mapsto b]}}
\and
\inferrule[T-CastE]{\tmj{\Gamma}{t}{E} \\ \elej{\Gamma}{\eta}{E}{E'}}{\tmj{\Gamma}{\ev{castE}\,t\,\eta}{E'}}
\and
\inferrule[T-Let]{\tmj{\Gamma}{t}{T} \\ \tmj{\Gamma, x : T}{u}{\weak{E}} \\ \clej{\Gamma, x : T}{f}{\uses{u}}{\weak{U}}}{\tmj{\Gamma}{\letc{x}{t}{u}{U}{f}}{E}}
\and
\inferrule[V-Lam]{\tmj{\mathit{body}(\Gamma, T)}{t}{E} \\ \clej{\mathit{body}(\Gamma, T)}{g}{\uses{t}}{\weak{A} \cup \cset{x}}}{\valj{\Gamma}{\lamcc{A}{\varrho}{\kappa, x}{T}{t}{g}}{\capt{(\pit[\kappa, x]{T}{E})}{A}}}
\and
\inferrule[V-Obj]{\mathit{obj}(\Gamma, T) \vdash^{\weak{A}}_{\!f} F}{\valj{\Gamma}{\litc{A}{W}{W^{c}}{F}}{T} \quad \text{where } T = \capt{(\obj{\ofLit{W, W^{c}}{\dom F}})}{A}}
\and
\inferrule[S-Cap]{\storej{\sigma}{\Gamma} \\ B \ne \scopeb}{\storej{\sigma, \kappa : B}{\Gamma, \kappa : B}}
\end{mathpar}
\color{black}
\caption{What \treeCC adds to Figures~\ref{fig:target-syntax} to~\ref{fig:target-capev} for scopes, and its changed rules. A binder $\kappa : (\instb{C})$ is written $\kappa \instb{C}$. Answer inclusion has two more rules, \rn{Ans-Pack} and \rn{Ans-Cong}, in Figure~\ref{fig:cc-pack}, and the dots in $\eta$ stand for their evidence. Section~\ref{sec:levels-scopes} defines the contexts $\mathit{scope}$, $\mathit{body}$ and $\mathit{obj}$. Here $\weak{X}$ is $X$ under the binders that the premise's context adds, in \rn{V-Obj} under $\varrho$ only. The other rules carry over, with answers for types.}
\label{fig:cc}
\end{figure}

Figure~\ref{fig:cc} gives what \treeCC adds for scopes. Its answers $E$ and their casts serve $\cfresh$ results (Section~\ref{sec:levels-fresh}), and until then an answer is a type. A capture set may contain the universal capability $\topc$, the compiler's top-level $\cany$ read as a constant. Besides star binders $\kappa : \cstar$, which stand for themselves, there are scope binders $\varrho : \scopeb$, the local $\cany$ of a scope, and instance binders $\kappa \instb{C}$, which stand for $C$. Scope and star binders are opaque. Four contexts open a scope:
\[
\begin{array}{@{}lcl@{}}
\mathit{scope}(\Gamma) & = & \Gamma, \varrho : \scopeb, \kappa : \cstar \\
\mathit{scope}(\Gamma, C) & = & \Gamma, \varrho : \scopeb, \kappa \instb{C} \\
\mathit{body}(\Gamma, T) & = & \Gamma, \varrho : \scopeb, \kappa : \cstar, x : T \\
\mathit{obj}(\Gamma, T) & = & \Gamma, \varrho : \scopeb, \self : T \mathrel{:=} \langle W, W^{c}, \dom F\rangle \qquad \text{for a literal } \litc{A}{W}{W^{c}}{F}
\end{array}
\]
In $\mathit{body}(\Gamma, T)$, $\varrho$ is the body's local $\cany$, $\kappa$ the arrow's capture binder, which a parameter written $\cany$ reads as, and $x$ the parameter. A call checks its argument $b$ at the domain with $\kappa$ at $\rt{b}$ (\rn{T-App}). So a parameter's $\cany$ is instantiated at the argument itself. The reference says it is instantiated at the capabilities passed in (\texttt{scoped-capabilities.md}, lines 32 to 34). The compiler creates a fresh capability per call, and the README of \treeCC lists this as a departure. Binding $\varrho$ first puts $\kappa$ and $x$ at the level of the body, so they may not flow out of it. A function shape $\pit[\kappa, x]{T}{E}$ binds $\kappa$ in $T$ and $E$ and binds $x$ in $E$, and a lambda $\lamcc{A}{\varrho}{\kappa, x}{T}{t}{g}$ binds all three in $t$ and $g$. In $\mathit{obj}(\Gamma, T)$, $\varrho$ is the local $\cany$ of a class, bound before the self. \rn{Sh-Pi} compares two domains in $\mathit{scope}(\Gamma)$, and a pack opens $\mathit{scope}(\Gamma, C)$.

A scope capability is $\topc$ or a scope binder. The level of a capture atom is the scope it belongs to, and $\topc$ is the outermost level:
\[
\begin{array}{@{}lcl@{}}
\lvl{\Gamma}{\topc} & = & \topc \\
\lvl{\Gamma}{\varrho} & = & \varrho \qquad \text{if } \Gamma \text{ binds } \varrho : \scopeb \\
\lvl{\Gamma}{y} & = & \varrho \qquad \text{if } y \text{ is another binder and } \varrho \text{ is the innermost scope binder left of } y \\
\lvl{\Gamma}{y} & = & \topc \qquad \text{if } y \text{ is another binder and no scope binder is left of } y \\
\lvl{\Gamma}{\sel{x}{\ell}} & = & \lvl{\Gamma}{x} \\
\lvlle{\Gamma}{c}{r} & \text{iff} & \lvl{\Gamma}{c} = \topc\text{, or } r \text{ is a capture binder and } \lvl{\Gamma}{c} \text{ is } r \text{ or left of } r \\
\lvlle{\Gamma}{C}{r} & \text{iff} & \lvlle{\Gamma}{c}{r} \text{ for every } c \in C
\end{array}
\]
The roots of Figure~\ref{fig:target-capev}, with clauses for $\topc$ and the two new bounds, are here called resolution, $\capsr{\Gamma}{n}{C}$. Roots apply expansion after resolution, and expansion opens each scope capability into $\topc$ and every opaque capture binder at or outside its level:
\[
\begin{array}{@{}lcl@{}}
\capsr{\Gamma}{n}{\topc} & = & \cset{\topc} \\
\capsr{\Gamma}{n}{\varrho} & = & \cset{\varrho} \qquad \text{if } \Gamma \text{ binds } \varrho : \scopeb \\
\capsr{\Gamma}{n}{\kappa} & = & \capsr{\Gamma'}{n}{C} \qquad \text{if } \Gamma = \Gamma', \kappa \instb{C}, \Gamma'' \\
\capsr{\Gamma}{}{C} & = & \textstyle\bigcup_{n} \capsr{\Gamma}{n}{C} \\
\expand{\Gamma}{r} & = & \cset{\topc} \cup \{\kappa \mid \kappa \text{ is opaque in } \Gamma \text{ and } \lvlle{\Gamma}{\kappa}{r}\} \qquad \text{if } r \text{ is a scope capability} \\
\expand{\Gamma}{c} & = & \cset{c} \qquad \text{for any other atom } c \\
\expand{\Gamma}{C} & = & \textstyle\bigcup_{c \in C} \expand{\Gamma}{c} \\
\rootsn{\Gamma}{n}{C} & = & \expand{\Gamma}{\capsr{\Gamma}{n}{C}}
\end{array}
\]
Expansion changes only scope capabilities, and over a store the only one is $\topc$ (Lemma~\ref{lem:lvl-rootfree}). So there the roots of a set are its resolution with $\topc$ opened into the platform, the sense of Section~\ref{sec:cc-types}. We read $\lvlle{\Gamma}{c}{r}$ as ``$c$ is at or outside the level of $r$''. So $\topc$ is universal only at the outermost level. Inside a scope the local $\varrho$ covers $\topc$ and every opaque capture binder at or outside its level. The level theorems speak of resolution, and effect safety of roots. Theorem~\ref{thm:target-uses} holds in \treeCC with these roots.

\begin{example}[The body context of the escape]\label{ex:lvl-escape}
Read \lstinline{File} as $\top$ and put the call at top level. The callback's body is typed in $\Gamma_{\!X4} = \mathit{body}(\cdot, \capt{\top}{\cset{\kappa_f}})$, the Lean's \lean{X4Ctx}. The returned closure uses $f$, so its annotation covers $f$, and the expected type puts that annotation below the $\cany$ outside the call, here $\topc$. So typing needs $\cset{f} \le \cset{\topc}$, which is argued only in comments of the Lean. $\Gamma_{\!X5}$ binds the scope binder last, the order this design rejects:
\[
\begin{array}{@{}l@{\qquad}l@{\qquad}l@{}}
\Gamma_{\!X4} = \varrho : \scopeb,\ \kappa_f : \cstar,\ f : \capt{\top}{\cset{\kappa_f}} & \lvl{}{\kappa_f} = \lvl{}{f} = \varrho & \expand{}{\topc} = \cset{\topc} \\
\Gamma_{\!X5} = \kappa_f : \cstar,\ f : \capt{\top}{\cset{\kappa_f}},\ \varrho : \scopeb & \lvl{}{\kappa_f} = \lvl{}{f} = \topc & \expand{}{\topc} = \cset{\topc, \kappa_f}
\end{array}
\]
In both contexts $\cset{f}$ resolves to $\cset{\kappa_f}$. In $\Gamma_{\!X4}$, $\kappa_f$ is not among the roots of $\cset{\topc}$, and in $\Gamma_{\!X5}$ it is. In $\Gamma_{\!X4}$ the roots of $\cset{\varrho}$ are $\cset{\topc, \varrho, \kappa_f}$: the callback's own $\cany$ covers its file. Status: the resolution in $\Gamma_{\!X4}$ is \lean{X4\_caps}, built by hand. The levels in $\Gamma_{\!X4}$ and of $f$ in $\Gamma_{\!X5}$ are decided, and the rest is on paper.
\end{example}

\subsection{The level rule and the changed rules}\label{sec:levels-rule}

Capture inclusion gains \rn{C-Level}. Its sides are singletons, and a larger set goes below $\cset{r}$ by $\ev{union}$. The rule lets an outer capability flow into an inner scope. It needs a scope capability on its right, so it never puts a capability below the arrow's binder, a star binder. In the context $\kappa_1 : \cstar, \varrho : \scopeb, \kappa_2 : \cstar$ it derives $\cset{\kappa_1} \le \cset{\varrho}$ and $\cset{\topc} \le \cset{\varrho}$, and its premise fails for $\varrho$ and $\kappa_2$ against $\topc$. In $\Gamma_{\!X5}$ the rule fires at $f$ and $\topc$, and in $\Gamma_{\!X4}$ it does not. Status: \lean{X2\_outer\_not\_inner}, in the Lean's \lean{X2Ctx}, and \lean{X5\_fires} are built by hand from decided premises, and \lean{X4\_no\_level} is decided.

\rn{CE-Inst} gives $\cset{\kappa} \equiv C$ for an instance binder $\kappa \instb{C}$, and $\ev{symm}$ and $\ev{eqToLe}$ give both inclusions. Five rules change for scopes: \rn{Sh-Pi}, \rn{V-Lam}, \rn{V-Obj}, \rn{T-App} and \rn{S-Cap}, which refuses a scope binder. \rn{T-Let} and \rn{K-Let} take a body at an answer.

\begin{figure}[t]
\footnotesize
\mathsrc
\[
\begin{array}{@{}lrcl@{}}
\text{Answers} & E & ::= & T \mid \exists(\kappa <:_c C).\,T \\
\text{Shapes} & S & ::= & \cdots \mid \forall(\kappa, x : T)\,E \\
\text{Terms} & t, u & ::= & \cdots \mid \klet \langle \kappa, x\rangle = t \kin u \\
\text{Contexts} & \Gamma & ::= & \cdots \mid \Gamma, \varrho : \scopeb \mid \Gamma, \kappa \instb{C}
\end{array}
\]
\begin{mathpar}
\inferrule[All]{\subj{\mathit{scope}(\Gamma)}{T_2}{T_1} \\ \subj{\mathit{body}(\Gamma, T_2)}{E_1}{E_2}}{\subj{\Gamma}{\forall(\kappa, x : T_1)\,E_1}{\forall(\kappa, x : T_2)\,E_2}}
\and
\inferrule[Sc-Level]{\Gamma \text{ binds } \varrho : \scopeb \\ \lvlle{\Gamma}{c}{\varrho}}{\subcj{\Gamma}{\cset{c}}{\cset{\varrho}}}
\and
\inferrule[Sc-Inst]{\Gamma \text{ binds } \kappa \instb{C}}{\subcj{\Gamma}{C}{\cset{\kappa}}}
\and
\inferrule[All-I]{\dtyuj{\weak{U} \cup \cset{x}}{\mathit{body}(\Gamma, T_1)}{t}{E} \\ T_1\ \mathit{wf}}{\dtyuj{\cempty}{\Gamma}{\lambda(x : T_1).\,t}{\capt{(\forall(\kappa, x : T_1)\,E)}{U}}}
\and
\inferrule[All-E]{\dtyuj{U}{\Gamma}{x}{\capt{(\forall(\kappa, z : T_1)\,E)}{C}} \\ \dtyuj{U}{\Gamma}{y}{T_1[\kappa \mapsto y]}}{\dtyuj{U}{\Gamma}{x\,y}{E[\kappa \mapsto y, z \mapsto y]}}
\and
\inferrule[\{\}-I]{\dtyuj{\weak{U} \cup \cset{x}}{\Gamma, \varrho : \scopeb, x : \capt{(\dmu{x}{S})}{U} \mathrel{:=} d}{d}{S} \\ \Distinct\,d}{\dtyuj{\cempty}{\Gamma}{\dnew{x}{d}}{\capt{(\dmu{x}{S})}{U}}}
\and
\inferrule[Sub]{\dtyuj{U}{\Gamma}{t}{E} \\ \subj{\Gamma}{E}{E'} \\ \subcj{\Gamma}{U}{U'}}{\dtyuj{U'}{\Gamma}{t}{E'}}
\and
\inferrule[Let-Ex]{\dtyuj{U_1}{\Gamma}{t}{\exists(\kappa <:_c C_0).\,T} \\ \subcj{\Gamma}{C_0}{U_2} \\ \dtyuj{\weak{U_2} \cup \cset{\kappa}}{\Gamma, \kappa : \cstar, x : T}{u}{\weak{E}}}{\dtyuj{U_1 \cup U_2}{\Gamma}{\klet \langle \kappa, x\rangle = t \kin u}{E}}
\end{mathpar}
\color{black}
\caption{The source of \treeCC beyond Figures~\ref{fig:source} and~\ref{fig:source-rules}. Answer inclusion $\subj{\Gamma}{E}{E'}$ is subtyping on types, with the rules \rn{Ans-Pack} and \rn{Ans-Cong} of Figure~\ref{fig:cc-pack} with the evidence erased, but not \rn{Ans-Trans}.}
\label{fig:cc-source}
\end{figure}

The source of \treeCC has the same scopes (Figure~\ref{fig:cc-source}), and the translation sends each binder to its target counterpart. Before typing, the source reads $\cany$ in the outer set of a parameter type as the arrow's capture binder. Elsewhere, where its \lean{AnyOk} allows it, it reads it as the enclosing scope binder, and at top level as the platform's set. It names no $\topc$, so \rn{Sc-Level} has a scope binder on the right. \rn{All}, \rn{All-I} and \rn{\{\}-I} open scopes as in the target, and Section~\ref{sec:levels-fresh} explains \rn{Sc-Inst} and \rn{Let-Ex}.

\subsection{Level inversion}\label{sec:levels-inversion}

Level inversion says that member-free evidence, however it is composed, never moves a capability outward. Member-free evidence, atoms and source derivations are defined by
\[
\begin{array}{@{}lcl@{}}
\mfree{g} & \text{iff} & g \text{ uses neither } \ev{member} \text{ nor } \ev{eqToLe}\text{, and } \mfree{a} \text{ for each atom } a \text{ under a } \ev{capvar} \text{ in } g \\
\mfree{a} & \text{iff} & \mfree{g} \text{ for the capture evidence } g \text{ of every cast and recap in } a \\
\mfree{\src{g}} & \text{iff} & \text{the source derivation } \src{g} \text{ uses none of } \rn{Sc-Inst}\text{, } \rn{Sc-Sel-Lower} \text{ and } \rn{Sc-Sel-Upper}
\end{array}
\]
The restriction is the lesson of the bad capture bounds of Example~\ref{ex:target-bad}. A lambda body may assume $\plec{\cset{\kappa}}{\cempty}$ through the object type of a binder, and only $\ev{member}$ and $\ev{eqToLe}$ read it.

\begin{theorem}[Level inversion, \lean{level\_inversion} and \lean{atom\_level\_inversion} in \treeCC and \treeCL]\label{thm:lvl-inv}
Let $\Gamma$ be a context and $r$ a capture atom.
\begin{enumerate}
\item If $\clej{\Gamma}{g}{C}{D}$, $\mfree{g}$ and $\lvlle{\Gamma}{\capsr{\Gamma}{}{D}}{r}$, then $\lvlle{\Gamma}{\capsr{\Gamma}{}{C}}{r}$.
\item If $\atomj{\Gamma}{a}{\capt{S}{C}}$, $\mfree{a}$ and $\lvlle{\Gamma}{\capsr{\Gamma}{}{C}}{r}$, then $\lvlle{\Gamma}{\capsr{\Gamma}{}{\cset{\rt{a}}}}{r}$.
\end{enumerate}
\end{theorem}

At $r = \topc$ it gives the last part of \lean{X2\_outer\_not\_inner}: no member-free evidence puts $\varrho$ or $\kappa_2$ below $\cset{\topc}$. It needs no store, and over a store it says nothing. A store opens no scope, so every atom there is at the outermost level and every level fact holds:

\begin{lemma}[A store opens no scope, \leanin{Store.Typed.rootFree}{\treeCC and \treeCL}]\label{lem:lvl-rootfree}
Let $\storej{\sigma}{\Gamma}$. Then $\Gamma$ binds no scope binder.
\end{lemma}

In a plain context the member-free restriction can go. A context is plain (\lean{Ctx.plain}) when $\ev{member}$ finds no capture proposition in it, and $\ev{eqToLe}$, \rn{CE-Def} and \rn{CE-Inst} find nothing to read:
\[
\begin{array}{@{}lcl@{}}
\mathit{plain}(\Gamma) & \text{iff} & \Gamma \text{ binds no instance binder, and every term binder of } \Gamma \text{ is opaque} \\
 & & \text{at a type } \capt{S}{C} \text{ with } \mathit{plain}(S) \\
\mathit{plain}(\bot) & = & \text{false} \\
\mathit{plain}(\sel{x}{A}) & = & \text{false} \\
\mathit{plain}(\pit[\kappa, x]{T}{E}) & = & \text{true} \\
\mathit{plain}(\boxt{T}) & = & \text{true} \\
\mathit{plain}(\obj{\tel}) & = & \text{true iff every proposition of } \tel \text{ is a bound } \pbnd{S'} \text{ with } \mathit{plain}(S')
\end{array}
\]

\begin{theorem}[Level inversion in a plain context, \leanin{level\_inversion\_plain}{\treeCC}]\label{thm:lvl-plain}
Let $\mathit{plain}(\Gamma)$ and let $r$ be a capture atom. If $\clej{\Gamma}{g}{C}{D}$ and $\lvlle{\Gamma}{\capsr{\Gamma}{}{D}}{r}$, then $\lvlle{\Gamma}{\capsr{\Gamma}{}{C}}{r}$.
\end{theorem}

\treeCL does not have this theorem. The escape at the target is its instance. No evidence at all puts $\cset{f}$ below $\cset{\topc}$ in $\Gamma_{\!X4}$ (\lean{X4\_no\_escape}, built by hand), since $\Gamma_{\!X4}$ is plain (\lean{X4Ctx\_plain}, decided) and $\cset{f}$ resolves to $\cset{\kappa_f}$, which is not at or outside $\topc$. The next theorem is about roots:

\begin{theorem}[Root inversion, \lean{root\_inversion} and \lean{atom\_root\_inversion} in \treeCC]\label{thm:root-inv}
Let $\Gamma$ be a context.
\begin{enumerate}
\item If $\clej{\Gamma}{g}{C}{D}$ and $\mfree{g}$, then $\caple{\Gamma}{C}{D}$.
\item If $\atomj{\Gamma}{a}{\capt{S}{C}}$ and $\mfree{a}$, then $\caple{\Gamma}{\cset{\rt{a}}}{C}$.
\end{enumerate}
\end{theorem}

It is Theorem~\ref{thm:target-canon} for member-free evidence with no store, the nearest Lean statement to Theorem~\ref{thm:mini-cap}. It separates two capabilities of one level, which Theorem~\ref{thm:lvl-inv} cannot (Example~\ref{ex:lvl-fresh}).

\begin{theorem}[Level inversion at the source, \leanin{source\_lvl\_safety}{\treeCC and \treeCL}]\label{thm:lvl-source}
Let $\src{\Gamma}$ be a well-formed source context, $\src{g}$ a derivation of $\subcj{\Gamma}{C}{D}$ with $\mfree{\src{g}}$, and $r$ a target capture atom. If $\lvlle{\trans{\Gamma}}{\capsr{\trans{\Gamma}}{}{\trans{D}}}{r}$, then $\lvlle{\trans{\Gamma}}{\capsr{\trans{\Gamma}}{}{\trans{C}}}{r}$.
\end{theorem}

The proof applies Theorem~\ref{thm:lvl-inv} to $\trans{g}$, which is typed and member-free. At the source the escape meets the platform, not $\topc$, since no translated set contains $\topc$ and a top-level $\cany$ reads as the platform's set. Take a hand-made $\src{\mathit{process} : \dall{f}{\capt{\mathit{File}}{\cset{\cany}}}{\top}}$ over the platform $\kappa_1, \kappa_2$. Its body context has no instance binder and no member to select, so every derivation there is member-free (\lean{Subcap.memberFree\_of\_plain}), and Theorem~\ref{thm:lvl-source} at $r = \topc$ shows that none proves $\cset{f} <:_c \cset{\kappa_1}$ (\lean{W5\_no\_escape}, built by hand).

\subsection{Running a body}\label{sec:levels-run}

\begin{figure}[t]
\footnotesize
\color{tgtcol}
\[
\begin{array}{@{}lcl@{}}
\ent{b} & = & [\varrho \mapsto \topc, \kappa \mapsto \rt{b}, z \mapsto b] \qquad \text{for the closure's binders } \varrho, \kappa, z \\
\appE{\omega}{\ev{plain}\,d} & = & \ev{cast}\,\omega\,d \\
\appE{\omega}{\ev{trans}\,\eta\,\eta'} & = & \appE{(\appE{\omega}{\eta})}{\eta'} \\
\appE{a}{\eta} & = & a \qquad \text{in every other case, and likewise for a value}
\end{array}
\]
\begin{mathpar}
\inferrule[K-Let]{\tmj{\Gamma, x : T}{u}{\weak{E}} \\ \clej{\Gamma, x : T}{f}{\uses{u}}{\weak{U}} \\ \contj{\Gamma}{K}{E}{T''}}{\contj{\Gamma}{K \triangleright \frlet{x}{u}{U}{f}}{T}{T''}}
\and
\inferrule[K-CastE]{\elej{\Gamma}{\eta}{E}{E'} \\ \contj{\Gamma}{K}{E'}{T''}}{\contj{\Gamma}{K \triangleright \ev{castE}\,\eta}{E}{T''}}
\end{mathpar}
\[
\begin{array}{@{}p{6.2em}@{}l@{}}
\rn{E-App} & \mst{\sigma}{K}{x\,b} \step \mst{\sigma}{K}{t\,\ent{b}} \qquad \text{if } \sigma(x) = \lamcc{A}{\varrho}{\kappa, z}{T}{t}{g} \\
\rn{E-AppCast} & \mst{\sigma}{K}{a\,b} \step \mst{\sigma}{K}{\ev{castE}\,(t\,\ent{b'})\,(\eta\,\ent{b})} \qquad \text{where } b' = \ev{cast}\,b\,(d\,\ent{b}) \\
 & \qquad \text{if } \sigma(\rt{a}) = \lamcc{A}{\varrho}{\kappa, z}{T}{t}{g},\ a \text{ is not a variable and } \chn{\sigma}{a}{n}{\form{pi}\,d\,\eta} \\
\rn{E-Proj} & \mst{\sigma}{K}{\ev{proj}\,a\,\ell\,h} \step \mst{\sigma}{K}{t[\varrho \mapsto \topc, \self \mapsto \rt{a}]} \\
 & \qquad \text{if } \sigma(\rt{a}) = \litc{A}{W}{W^{c}}{F} \text{ and } F \ni \ell = t\,\ann{g} \\
\rn{E-CastEPush} & \mst{\sigma}{K}{\ev{castE}\,t\,\eta} \step \mst{\sigma}{K \triangleright \ev{castE}\,\eta}{t} \\
\rn{E-CastEVal} & \mst{\sigma}{K \triangleright \ev{castE}\,\eta}{v} \step \mst{\sigma}{K}{\appE{v}{\eta}} \\
\rn{E-CastEAtom} & \mst{\sigma}{K \triangleright \ev{castE}\,\eta}{a} \step \mst{\sigma}{K}{\appE{a}{\eta}}
\end{array}
\]
\color{black}
\caption{The machine of \treeCC where it differs from Figure~\ref{fig:target-machine}, with the answer cast $\appE{\omega}{\eta}$ of an atom or value $\omega$. $\contj{\Gamma}{K}{E}{T''}$ says that $K$ accepts an answer of type $E$ and produces a result of type $T''$. The first equation that applies is used, and the last covers the cases no typed state reaches. \rn{E-AppCastRefl} substitutes $\ent{b}$ as \rn{E-App} does.}
\label{fig:cc-machine}
\end{figure}

\rn{E-App} sends $\varrho$ to $\topc$, $\kappa$ to $\rt{b}$ and the parameter to $b$ (Figure~\ref{fig:cc-machine}). A running program is the outermost scope, so the local $\cany$ of a body it enters becomes the universal one:

\begin{lemma}[Entering a body, \lean{Subst.Typed.enter} and the \lean{HasType.subst} lemmas in \treeCC and \treeCL]\label{lem:lvl-enter}
Let $\Gamma$ bind no scope binder, and let $\atomj{\Gamma}{b}{T[\kappa \mapsto \rt{b}]}$. Then every typing judgment that holds in $\mathit{body}(\Gamma, T)$ holds in $\Gamma$ after the substitution $\ent{b}$ of Figure~\ref{fig:cc-machine}.
\end{lemma}

Level facts about $\varrho$ become level facts about $\topc$, which hold since $\Gamma$ binds no scope binder. The machine applies the lemma over its store (Lemma~\ref{lem:lvl-rootfree}). Store extension appends no opaque capture binder, which would grow the expansion of $\topc$.

In \treeCC, Theorem~\ref{thm:target-canon} gains the cases \rn{C-Level} and \rn{CE-Inst}. The \rn{C-Level} case needs no store: what is at or outside the level of $r$ has its roots in the expansion of $r$. An instance binder has the roots of its set, which gives \rn{CE-Inst}.

\treeCC transcribes the running example of Section~\ref{sec:mini} and decides its typings. Its closures never mention their scope binders, so \lean{C6\_run} takes $\mathit{prog}_1$ through the seven steps of Table~\ref{tab:mini-run}, and \lean{C6\_safe} is effect safety for $\mathit{prog}_2$ (both built by hand).

\subsection{Fresh results}\label{sec:levels-fresh}

\begin{figure}[t]
\footnotesize
\color{tgtcol}
\[
\begin{array}{@{}l@{}}
\begin{array}{@{}l@{\quad}rcl@{}}
\text{Packed atoms} & \bar{a} & ::= & a \mid \ev{pack}\,C\,g\,d\,a \\
\text{Values} & v & ::= & \cdots \mid \ev{pack}\,C\,g\,d\,v \\
\text{Terms} & t, u & ::= & \bar{a} \mid \cdots \mid \letex{\kappa}{x}{t}{u}{U}{g}{f} \\
\text{Answer evidence} & \eta & ::= & \cdots \mid \ev{pack}\,C\,g\,d \mid \ev{cong}\,g\,d \\
\text{Continuations} & K & ::= & \cdots \mid K \triangleright L \qquad \text{where } L = \frletex{\kappa}{x}{u}{U}{g}{f}
\end{array} \\[6pt]
\begin{array}{@{}lcl@{\qquad}lcl@{}}
\rt{\ev{pack}\,C\,g\,d\,a} & = & \rt{a} & \uses{\letex{\kappa}{x}{t}{u}{U}{g}{f}} & = & \uses{t} \cup U \\
\uses{\bar{a}} & = & \cset{\rt{\bar{a}}} & \uses{K \triangleright L} & = & \uses{K} \cup U
\end{array} \\[3pt]
\begin{array}{@{}lcl@{}}
\appE{\omega}{\ev{pack}\,C\,g\,d} & = & \ev{pack}\,C\,g\,d\,\omega \\
\ev{capt}\,e\,f \mathbin{;} \ev{capt}\,e'\,f' & = & \ev{capt}\,(\ev{trans}\,e\,e')\,(\ev{trans}\,f\,f') \\
\appE{(\ev{pack}\,C\,g\,d\,\omega)}{\ev{cong}\,g'\,d'} & = & \ev{pack}\,C\,(\ev{trans}\,g\,g')\,(d \mathbin{;} d')\,\omega
\end{array}
\end{array}
\]
\begin{mathpar}
\inferrule[Ans-Pack]{\clej{\Gamma}{g}{C}{C_0} \\ \lej{\mathit{scope}(\Gamma, C)}{d}{\weak{T'}}{T}}{\elej{\Gamma}{\ev{pack}\,C\,g\,d}{T'}{\exc{\kappa}{C_0}{T}}}
\and
\inferrule[Ans-Cong]{\clej{\Gamma}{g}{C_0}{C_0'} \\ \lej{\mathit{scope}(\Gamma)}{d}{T}{T'}}{\elej{\Gamma}{\ev{cong}\,g\,d}{\exc{\kappa}{C_0}{T}}{\exc{\kappa}{C_0'}{T'}}}
\and
\inferrule[T-PackA]{\atomj{\Gamma}{a}{T'} \\ \clej{\Gamma}{g}{C}{C_0} \\ \lej{\mathit{scope}(\Gamma, C)}{d}{\weak{T'}}{T}}{\tmj{\Gamma}{\ev{pack}\,C\,g\,d\,a}{\exc{\kappa}{C_0}{T}}}
\and
\inferrule[T-PackV]{\valj{\Gamma}{v}{T'} \\ \clej{\Gamma}{g}{C}{C_0} \\ \lej{\mathit{scope}(\Gamma, C)}{d}{\weak{T'}}{T}}{\tmj{\Gamma}{\ev{pack}\,C\,g\,d\,v}{\exc{\kappa}{C_0}{T}}}
\and
\inferrule[T-LetEx]{\tmj{\Gamma}{t}{\exc{\kappa}{C_0}{T}} \\ \clej{\Gamma}{g}{C_0}{U} \\ \tmj{\Gamma, \kappa : \cstar, x : T}{u}{\weak{E}} \\ \clej{\Gamma, \kappa : \cstar, x : T}{f}{\uses{u}}{\weak{U} \cup \cset{\kappa}}}{\tmj{\Gamma}{\letex{\kappa}{x}{t}{u}{U}{g}{f}}{E}}
\and
\inferrule[K-LetEx]{\clej{\Gamma}{g}{C_0}{U} \\ \tmj{\Gamma, \kappa : \cstar, x : T}{u}{\weak{E}} \\ \clej{\Gamma, \kappa : \cstar, x : T}{f}{\uses{u}}{\weak{U} \cup \cset{\kappa}} \\ \contj{\Gamma}{K}{E}{T''}}{\contj{\Gamma}{K \triangleright L}{\exc{\kappa}{C_0}{T}}{T''}}
\end{mathpar}
\[
\begin{array}{@{}p{6.2em}@{}l@{}}
\rn{E-LetEx} & \mst{\sigma}{K}{\letex{\kappa}{x}{t}{u}{U}{g}{f}} \step \mst{\sigma}{K \triangleright L}{t} \\
\rn{E-Unpack} & \mst{\sigma}{K \triangleright L}{\ev{pack}\,C\,g_0\,d\,a} \step \mst{(\sigma, \kappa \instb{C})}{K}{u[x \mapsto \ev{cast}\,a\,d']} \\
 & \qquad \text{where } d' = d[\varrho' \mapsto \topc, \kappa' \mapsto \kappa] \text{ for the two binders } \varrho', \kappa' \text{ that } \mathit{scope}(\Gamma, C) \text{ adds} \\
\rn{E-UnpackV} & \mst{\sigma}{K \triangleright L}{\ev{pack}\,C\,g_0\,d\,v} \step \mst{(\sigma, \kappa \instb{C}, x = \mathit{core}(v'))}{K}{\mathit{adjust}(u, v')} \\
 & \qquad \text{where } v' = \ev{cast}\,v\,d' \text{ and } d' \text{ is as for \rn{E-Unpack}}
\end{array}
\]
\color{black}
\caption{Packing and unpacking in \treeCC, and the steps that unpack. An atom term, here and in \rn{E-CastEAtom} and the last clause of Figure~\ref{fig:cc-machine}, is now a packed atom $\bar{a}$, an atom or a pack of one. No call, projection or unboxing applies to a pack. The two clauses of $\appE{\omega}{\eta}$ come before that last one.}
\label{fig:cc-pack}
\end{figure}

A result type written with $\cfresh$ promises a new capability at each call. \treeCC reads it as an existential answer $\exc{\kappa}{C_0}{T}$, as System Capybara also does \cite{xu2026capybara}: there is a $\kappa$ below $C_0$ such that the result has type $T$. Answers stand only as codomains of function shapes and as types of terms. Packing (Figure~\ref{fig:cc-pack}) is an answer inclusion, \rn{Ans-Pack}, and a wrapper on atoms and values, \rn{T-PackA} and \rn{T-PackV}, which unpacking matches. No rule types a packed value at a plain type, so a packed value is never stored. The bound $C_0$ is what the unpacking charges for its witness, by $g : C_0 \le U$ in \rn{T-LetEx}. Without it the unpacking would charge the enclosing scope's $\cany$, which at top level is $\topc$, and effect safety would then say nothing. At run time the unpacking stores the witness as an instance binder $\kappa \instb{C}$, which \rn{CE-Inst} reads. The source reads $\cfresh$ the same way (Figure~\ref{fig:cc-source}). \rn{Sub} packs by answer inclusion, and \rn{Let-Ex} unpacks. A pack needs its witness below the opened binder, so \rn{Sc-Inst} has one direction.

The theorems of Section~\ref{sec:source-thms} take a program typed at a type, not an answer, as \rn{K-Nil} accepts only types. \lean{Platform.root\_iff} of \treeCC needs $\topc \notin \trans{U}$, which translated sets meet.

\begin{remark}[Existentials are isolated, \leanin{no\_ex\_le\_ty}{\treeCC and \treeCL}]\label{rem:lvl-ex}
For every $\Gamma$, $\eta$, $C_0$, $T$ and $T'$, there is no derivation of $\elej{\Gamma}{\eta}{\exc{\kappa}{C_0}{T}}{T'}$.
\end{remark}

It is immediate from the rules. No rule of answer inclusion has an existential on the left and a plain type on the right, and \rn{Ans-Trans} keeps an existential an existential (\lean{ex\_stays\_ex}).

\begin{example}[\textup{\texttt{withFile}} with a fresh result]\label{ex:lvl-fresh}
Let the callback of Section~\ref{sec:levels-program} hide its file behind a $\cfresh$ result. With \lstinline{File} read as $\top$, it has the type
\[
\mathit{op} : \capt{(\pit[\kappa, z]{\capt{\top}{\cset{\kappa}}}{\exc{\kappa'}{\cset{z}}{\capt{(\pit[\kappa'', w]{\Unit}{\Unit})}{\cset{\kappa'}}}})}{\cempty} .
\] The body of \lstinline{withFile} is typed in $\mathit{body}(\mathit{op} : \ldots, \capt{\top}{\cset{\kappa_f}})$, which binds $\varrho_b$, $\kappa_f$ and $f$. It calls $\mathit{op}$ on its file. By Remark~\ref{rem:lvl-ex}, no evidence casts the answer to a plain type, so the body unpacks it by a $\kw{let}\,\langle \kappa', x\rangle$ that declares $\cset{f}$ and charges it by $\ev{refl}\,\cset{f}$ (\rn{T-LetEx}). It opens $\kappa' : \cstar$ without its bound, as a new capability, and $x$ at the set $\cset{\kappa'}$. The body is typed, and $\kappa'$ has the level $\varrho_b$. By Theorem~\ref{thm:root-inv}, no member-free evidence puts $\cset{\kappa'}$ below $\cset{f}$. The context is plain, so by Theorem~\ref{thm:lvl-plain} no evidence at all puts $\cset{x}$ below $\cset{\topc}$. So $\cfresh$ hides $f$ behind $\kappa'$. Status: the typing and the level are decided, and \lean{Y4\_isolation} and \lean{Y4\_no\_escape} are built by hand.
\end{example}

\subsection{What is not proved}\label{sec:levels-open}

No theorem of \treeCC says that the escape program has no derivation, in the source or in the target. \lean{X4\_no\_escape} answers one subcapturing question in one hand-made target context, and no translated program poses it. \lean{W5\_no\_escape} answers one in one hand-made source context, at $\cset{\kappa_1}$, not at the platform's set $\cset{\kappa_1, \kappa_2}$. That typing the escape needs $\cset{f} \le \cset{\topc}$ (Example~\ref{ex:lvl-escape}) and that \rn{C-Level} matches the compiler's own level check are said only in comments of the Lean. Whole-program rejection is the verdict of the front end (Section~\ref{sec:fe-reject}), and its certificate rules out member-free derivations of one goal. The use after close of the reference is not modelled, since nothing closes a file.

\rn{Ans-Pack} puts no level condition on its witness, so a local $\cany$ may flow into a result $\cfresh$, which the reference forbids. Two calls of a function with a $\cfresh$ result open two star binders, and in that typing context no capture evidence puts either below the other (\lean{two\_calls\_incomparable}, built by hand). At run time both witnesses become instance binders at the same set, and \rn{CE-Inst} relates them (on paper). So the two results are incomparable only while the program is checked.
